%%%%%%%%%%%%%%%%%%%%%%%%%%%%%%%%%%%%%%%%%%%%%%%%%%%%%%%%%%%%%%%%%%%%%%%%%%%%%%%%%
%	TASK 2 - Algorithm Development
%   Maps to brief Tasks 2.1 (GA) and 2.2 (Constraint Handling via Repair).
%%%%%%%%%%%%%%%%%%%%%%%%%%%%%%%%%%%%%%%%%%%%%%%%%%%%%%%%%%%%%%%%%%%%%%%%%%%%%%%%%
\lhead{Task 2. \emph{Algorithm Development}}
\section{Task 2: Algorithm Development}
\label{Task2}

This chapter describes the genetic algorithm used throughout the project and
the repair mechanism that handles the permutation constraint. Section
\ref{T2:GADesign} defines the representation and the five evolutionary
operators. Section \ref{T2:Repair} explains why the primary recombination
operator produces infeasible offspring on every generation and describes the
repair step that converts those offspring back into valid Hamiltonian tours.
The empirical comparison of repair against alternative constraint-handling
approaches is reported in Section \ref{T3:RepairCritique}.

\subsection{Genetic Algorithm Design (Task 2.1)}
\label{T2:GADesign}

\subsubsection{Solution Representation}
\label{T2:Representation}
Each candidate solution is encoded as a permutation of the city index set
$\{0, 1, \ldots, n-1\}$. A chromosome is an ordered tuple where the
$i$-th gene indicates the $i$-th city visited. The tour is closed by returning
to the first gene. For the eight-city illustration used later, the
chromosome $[3,0,7,1,5,2,6,4]$ denotes the tour
$3 \to 0 \to 7 \to 1 \to 5 \to 2 \to 6 \to 4 \to 3$. The fitness of a
chromosome $\pi$ is its closed Euclidean tour length,
\begin{equation}
f(\pi) \;=\; \sum_{i=0}^{n-1} d_{\pi_i,\,\pi_{(i+1) \bmod n}},
\end{equation}
where $d_{ij}$ is the Euclidean distance between cities $i$ and $j$.
Because fitness is computed by indexing the distance matrix at the gene
values, the calculation is only meaningful when the chromosome is a valid
permutation. A chromosome with duplicate or missing cities still produces a
number, but that number is not the length of a Hamiltonian tour and should
not be allowed into selection (see Section \ref{T2:Infeasibility}).

\subsubsection{Population Initialisation}
\label{T2:Init}
The initial population is generated as $N$ independent random permutations
of the city set, produced with
\texttt{numpy.random.default\_rng(seed).permutation(n)}. Each chromosome is
sampled independently, so the initial population is drawn uniformly from
the space of $n!$ permutations and starts with high diversity. No
heuristic seeding (such as a nearest-neighbour tour) is used, so
the convergence behaviour reported in Section \ref{T4:Convergence}
can be attributed to the evolutionary operators rather than to a
pre-seeded population. A single seed controls every random component of the GA,
making each run reproducible and allowing paired-seed comparisons in
Section \ref{T3:Sensitivity}.

\subsubsection{Selection}
\label{T2:Selection}
Two selection methods are implemented and controlled through a single
configuration parameter, so the choice of parent-selection rule can be
varied in the same way as any numerical hyperparameter.

Tournament selection is the primary operator. Given a population $P$ of
size $N$, $k$ candidates are drawn uniformly at random without replacement
and the candidate with the lowest tour length is returned,
\begin{equation}
s \;=\; \arg\min_{i \in C} f(i), \qquad C \subseteq P,\; |C| = k.
\end{equation}
The tournament size $k$ controls selection pressure: small $k$ gives weaker
individuals a better chance, while large $k$ narrows the gene pool quickly.
The primary configuration uses $k = 3$, which keeps enough randomness to
delay premature convergence while still favouring the better half of the
population. The sweep in Section \ref{T3:Sensitivity} varies $k$
explicitly.

Roulette wheel selection is implemented as an alternative operator for
comparison. Because the TSP is a minimisation problem, raw fitness cannot
be used as a selection weight directly. Instead, each individual is
assigned the inverse fitness and the weights are normalised to a
probability distribution,
\begin{equation}
p_i \;=\; \frac{1/f_i}{\sum_{j=1}^{N} 1/f_j}.
\end{equation}
The roulette draw then picks a parent with probability $p_i$. A single
\texttt{select} function checks the \texttt{selection\_method} field in the
GA configuration and calls the appropriate method, so the main loop
is the same for both selection rules.

\subsubsection{Crossover}
\label{T2:Crossover}
The primary recombination operator is naive single-point crossover.
Given two parents $A$ and $B$ of length $n$ and a cut point
$pt \in \{1, \ldots, n-1\}$, the child is constructed by concatenation,
\begin{equation}
\mathrm{child} \;=\; A[\,0\,{:}\,pt\,] \,\Vert\, B[\,pt\,{:}\,n\,].
\end{equation}
This operator copies positions independently, without checking whether the
inserted gene values have already appeared elsewhere in the child. It does
not preserve the permutation property: the same city may appear twice and
another may disappear. This choice is intentional. Pairing naive crossover
with the repair step means the repair mechanism must do real work on every
generation, not just handle rare edge cases. This setup matches the
assignment's focus on repair as the main constraint-handling contribution.

Three feasibility-preserving operators are also implemented as
alternatives for the comparison in Section \ref{T3:RepairCritique}.
Partially Mapped Crossover (PMX) \cite{goldberg1985alleles} copies a
segment $A[p_1{:}p_2]$ into the child and follows the segment mapping for
each displaced value from $B$ until it finds an empty position; remaining
slots are filled directly from $B$. Order Crossover (OX),
from \cite{davis1985applying} and refined in \cite{oliver1987study}, also
copies the segment from $A$ and then walks $B$ in wrap-around order
starting after $p_2$, skipping any value already in the segment. Cycle
Crossover (CX) finds position-wise cycles between the two parents and
takes whole cycles from each parent in turn, so every position is
filled once and every city appears once. The reasoning behind PMX, OX,
and CX, along with comparative reviews of permutation operators, is
covered in \cite{larranaga1999genetic} and
\cite{cicirello2023permutationoperators}. The worked example in Section
\ref{T2:Infeasibility} shows all three operators producing valid
permutations on the same parents that break naive crossover.

\subsubsection{Mutation}
\label{T2:Mutation}
Swap mutation is the sole mutation operator used in the primary pipeline.
On each call, a random number is drawn from $[0,1)$. If it falls
below the mutation rate $r_m$, two distinct gene positions $i$ and $j$ are
picked without replacement and the gene values at those positions are
swapped. The chromosome length, the set of cities, and the permutation
property are all preserved by design, so swap mutation does not introduce
any infeasibility. The change is local: exactly two edges in
the tour change, which makes it suitable for late-stage refinement.
Inversion mutation, where the segment between two cut points is
reversed in place, is implemented in the code base as an alternative
operator but is not used in any of the primary experiments in this report.

\subsubsection{Main Loop}
The five operators are combined into a generational GA with elitism.
At each generation the top $e$ individuals are copied unchanged into the
next population. The remaining slots are filled by repeated selection,
crossover (with probability $r_c$), mutation, and (if enabled) repair.
The elitism count is fixed at $e = 2$ throughout the primary experiments,
which keeps the current best solution safe while leaving most of the
population open to recombination. Algorithm \ref{alg:ga} shows the loop
in detail.

\begin{algorithm}[h]
\caption{Generational GA with Optional Repair}
\label{alg:ga}
\begin{algorithmic}[1]
\State Initialise population $P$ of $N$ random permutations
\State Evaluate fitness $f(p)$ for all $p \in P$
\For{$g = 1$ to $G_{\max}$}
    \State $P' \gets$ top-$e$ elites from $P$ \Comment{elitism}
    \While{$|P'| < N$}
        \State Select parents $p_1, p_2$ from $P$ via tournament or roulette
        \If{$\mathrm{rand}() < r_c$}
            \State $c \gets \mathrm{Crossover}(p_1, p_2)$ \Comment{naive single-point in the primary pipeline}
        \Else
            \State $c \gets \mathrm{copy}(p_1)$
        \EndIf
        \State $c \gets \mathrm{Mutate}(c, r_m)$ \Comment{swap mutation}
        \If{repair enabled}
            \State $c \gets \mathrm{Repair}(c)$
        \EndIf
        \State Append $c$ to $P'$
    \EndWhile
    \State $P \gets P'$; evaluate fitness for all $p \in P$
\EndFor
\State \Return best individual in $P$
\end{algorithmic}
\end{algorithm}

\subsection{Constraint Handling via Repair Mechanism (Task 2.2)}
\label{T2:Repair}

The TSP imposes a single hard constraint on the chromosome: every city
must appear exactly once. This constraint cannot be relaxed without
changing the problem, so the GA either uses recombination operators that
preserve it by design or repairs violations after they occur.
Surveys of evolutionary constraint handling describe the same three
options --- preservation, repair, and penalty
\cite{coello2002constraint,michalewicz1996evolutionary}. The primary
pipeline of this project uses the repair approach, paired with naive
single-point crossover so that the repair step does real work on every
offspring.

\subsubsection{Why Infeasibility Arises}
\label{T2:Infeasibility}
Naive single-point crossover treats the chromosome like an ordinary
fixed-position string. It copies gene positions independently and does
not check whether the inserted value already appears elsewhere in the
child. In a permutation representation the constraint applies to the
whole chromosome, not to individual positions, so this position-by-position
copy is what breaks the city set.

Consider the following worked example. The parents are
\begin{align*}
A &= (0,1,2,3,4,5,6,7) \\
B &= (2,6,4,0,5,7,1,3)
\end{align*}

and the single cut point is $pt = 4$. The child constructed by naive
single-point crossover is

\begin{align*}
\mathrm{child} &= A[\,0{:}4\,] \,\Vert\, B[\,4{:}8\,] \\
&= (0,1,2,3,5,7,1,3)
\end{align*}

City $1$ now appears at positions $1$ and $6$, and city $3$ appears at
positions $3$ and $7$. Cities $4$ and $6$ never appear. The chromosome
has the correct length but is not a Hamiltonian tour. The same
problem occurs on kroA100, where all thirty naive-crossover offspring
tested were infeasible, with twenty-four duplicates on average.

The real problem is the effect on selection. Fitness evaluation
on $(0,1,2,3,5,7,1,3)$ still returns a number, but that number sums
distances over an edge set that visits city $1$ twice, visits city $3$
twice, and skips cities $4$ and $6$ entirely. The result tends to be
artificially small because cheap edges are counted repeatedly while
expensive missing-city edges are left out. Selection compares this
number against the actual tour lengths of feasible competitors as if
they were comparable. Tournament and roulette selection then favour the
infeasible chromosome, its invalid genetic material spreads through
recombination, and within a few generations the population moves towards
short pseudo-tours rather than short tours. The infeasibility analysis in
\cite{cicirello2023permutationoperators} traces the same failure pattern
across the permutation operator family.

Comparing with the permutation-preserving operators makes this
clear. On the same parents and the same cut points:

\begin{itemize}
	\item PMX returns $(5,6,2,3,4,7,1,0)$.
	\item OX returns $(0,5,2,3,4,7,1,6)$.
	\item CX returns $(0,6,2,3,4,5,1,7)$
\end{itemize}


All three are valid permutations because
each operator has a feasibility-preserving rule built into the copy step
itself. The repair mechanism described below achieves the same
result by acting after a naive copy rather than during one.

\subsubsection{Repair Design}
\label{T2:RepairDesign}
The repair function has two phases: diagnose the infeasible chromosome,
then correct it. The diagnosis phase does a single left-to-right scan.
A \emph{seen} set records the first occurrence of each city in scan
order. Any later occurrence of the same city is marked as a duplicate
position. After the scan, the missing set is computed as
$\{0, 1, \ldots, n-1\} \setminus \mathrm{seen}$. Because the chromosome
length is fixed at $n$, the number of duplicate positions equals the
number of missing cities, so every duplicate can be replaced with a
unique missing city.

The correction phase assigns the missing cities to the duplicate
positions. The order of assignment differs across three strategies:
\begin{enumerate}[nosep]
    \item \textbf{Random insertion} (the primary strategy) shuffles the
          missing cities with the GA's seeded random generator and
          places them at the duplicate positions in shuffled order.
    \item \textbf{Position insertion} places the missing cities in
          sorted order at the duplicate positions in left-to-right
          order. The assignment is deterministic, so the same
          infeasible chromosome always repairs to the same tour.
    \item \textbf{Nearest-neighbour insertion} processes the duplicate
          positions left to right and, at each position, chooses the
          still-unassigned missing city whose Euclidean distance to the
          immediately preceding gene in the chromosome is smallest.
\end{enumerate}
Random insertion is the primary choice. It does not add any bias from the
problem geometry, so the repair step only provides feasibility without
guiding optimisation. The GA's selection pressure remains the only
driver of solution quality. Random insertion also tends to keep
population diversity higher than the deterministic alternatives,
because two infeasible chromosomes with the same duplicates and missing
cities can still be repaired into different tours. The
feasibility-preserving repair literature reports the same trade-off
between feasibility and diversity
\cite{padhye2015feasibilitypreserving}. Algorithm \ref{alg:ga} places
the repair call after mutation, so any chromosome reaching the fitness
evaluation step at the end of the loop is guaranteed to be a valid
permutation.

The repair function is shown step by step below. On the eight-city
chromosome $(3,1,4,1,5,2,6,3)$, the diagnosis phase finds duplicate
positions $\{3, 7\}$ and missing cities $\{0, 7\}$. Random insertion may
return $(3,1,4,7,5,2,6,0)$; position insertion always returns
$(3,1,4,0,5,2,6,7)$; nearest-neighbour insertion on the toy octagon
returns $(3,1,4,7,5,2,6,0)$. All three are valid permutations. A
property test runs the repair on five hundred random corruptions across
all three strategies and produces a valid permutation in every case.

\subsubsection{Effect on Quality, Diversity, and Convergence}
\label{T2:RepairEffect}
The repair step has three intended effects on the GA's behaviour.
First, it ensures fitness validity: every chromosome that reaches the
evaluation step in Algorithm \ref{alg:ga} is a valid permutation, so
the tour length returned by the distance matrix is the actual length of
a Hamiltonian tour. Without repair, the fitness function still returns a
number, but it is distorted by duplicate and missing-city edges and
cannot be compared against valid tours. Second, random insertion adds
variation at the duplicate positions, which helps maintain population
diversity. Third, because repair keeps the search inside the feasible
region, the convergence curve reflects actual tour-length improvement
rather than worsening constraint violations. Section
\ref{T3:RepairCritique} confirms these effects quantitatively.
